\section*{Summary}
\vspace{-2ex}\titlerule\vspace{4ex}

Accurate temperature measurement is critical in industrial metal forming and heat-treatment processes, where pyrometers provide non-contact sensing at extreme temperatures. However, raw pyrometer signals suffer from two key problems that are addressed in this individual thesis contribution: systematic calibration errors caused by the temperature-dependent emissivity of metal surfaces, and the large volume of raw data generated during continuous production runs that must be stored and transmitted efficiently.

This thesis addresses these two problems through the design, implementation, and evaluation of automated sensor calibration and data compression stages within a modular pre-processing pipeline, structured around two acceptance test procedures. In the calibration stage~(ATP-2), four classical methods were implemented and compared: mean offset correction, linear regression~(RMSE\,=\,168.3\,\textdegree C), polynomial regression~(degree\,=\,2), and piecewise linear calibration~(3 segments). These were evaluated against four machine learning approaches: Random Forest regression~(100 trees, RMSE\,=\,100.1\,\textdegree C), MLP neural network~(architecture 64--128--64, RMSE\,=\,103.6\,\textdegree C), Gradient Boosting~(200 trees, RMSE\,=\,102.8\,\textdegree C), and Support Vector Regression~(RBF kernel, RMSE\,=\,107.8\,\textdegree C). Testing on the NIST open-source additive manufacturing dataset~(Layer~01) showed that all four machine learning methods outperformed all classical regression methods by more than 35\,\%, with Random Forest achieving the best result. The ATP-2 target of~$R^2 > 0.99$ was satisfied by all four ML methods.

In the compression stage~(ATP-3), three additional compression methods were implemented and evaluated as an individual contribution: Delta Encoding~(int16, near-lossless, RMSE\,$\approx$\,0.04\,\textdegree C, ratio\,$\approx$\,2$\times$), Variational Autoencoder~(VAE, bottleneck\,=\,3, ratio\,=\,10.7$\times$), and Deep Autoencoder~(bottleneck\,=\,3, ratio\,=\,10.7$\times$). Delta Encoding achieved near-lossless compression on the slowly varying calibrated pyrometer signal. The Deep Autoencoder achieved RMSE\,=\,113.0\,\textdegree C at~10.7$\times$, outperforming the standard Autoencoder from the parallel thesis by~33.8\,\%.

The individual contribution of this thesis is the systematic investigation of sensor calibration methods and additional data compression methods for pyrometer time-series, demonstrating that machine learning calibration methods substantially outperform classical regression approaches and that Delta Encoding provides near-lossless compression for slowly varying industrial temperature signals. The pipeline is delivered as a modular, open-source Python implementation suitable for integration into industrial monitoring systems.
